\begin{abstract}
Although Large Language Models (LLMs) exhibit strong inherent fault tolerance and mask the effects of most GPU soft errors, a small fraction of these errors can still cause severe semantic deviations in their outputs, leaving LLM inference vulnerable to Silent Data Corruption (SDC) with potentially serious consequences for safety-critical applications such as autonomous driving and robotics.
Detecting these consequential faults is challenging because the magnitude of numerical anomalies alone does not reliably indicate their semantic impact.
Our key observation is that features in adjacent LLM layers exhibit strong predictability during fault-free execution, while Significant SDCs tend to disrupt this relationship, increasing the discrepancy between a layer's predicted and observed features.
Building on this observation, we propose SIEVE, a lightweight framework for detecting Significant SDCs through deviations from learned inter-layer predictions.
SIEVE trains a Nominal Inter-Layer Predictor on clean execution traces to predict the current layer's projected self-attention features from those of the preceding layer.
At inference time, it extracts statistical and geometric discrepancies between the predicted and observed features of the current layer, aggregates them across monitored layer pairs and decoding steps into a Discrepancy Profile, and uses a lightweight XGBoost classifier to identify Significant SDCs.
Across three LLMs and three VQA benchmarks, SIEVE achieves 96.38\% precision, 91.49\% recall, and 93.78\% F1 on the primary evaluation set, with a false positive rate of 0.11\%.
Compared with representative SDC detection baselines, SIEVE improves F1 by 11.38 percentage points while introducing an average end-to-end overhead of 6.65\%.
\end{abstract}

